\section{Tire Relaxation, Grip States \& Advanced Dynamics}

\label{sec:tire-relaxation}

This section covers tire force relaxation (the lag between instantaneous slip and realized force), grip state transitions, and per-wheel grip smoothing. Proper handling of tire relaxation is essential for realistic drifting and sliding behavior.

\subsection{Tire Relaxation Length}

The relaxation distance $\lambda$ represents how far the tire must roll before reaching steady-state force after a sudden slip change:

\begin{equation}
  \lambda = \text{relaxationLength} \approx 0.3 \text{ m}
  \label{eq:relaxation-length}
\end{equation}

This converts to a time constant based on longitudinal speed:

\begin{equation}
  \tau(v) = \frac{\lambda}{\max(|v_{\text{long}}|, 0.5 \text{ m/s})}
  \label{eq:relaxation-tau}
\end{equation}

\textbf{Physical Meaning}: Tire relaxation models the transient response of tire slip. The relaxation length is a material property; at higher speeds, the relaxation time (distance/speed) is shorter, causing faster force response.

\textbf{Source Code}: \coderef{tires.js}{332--343}

\begin{lstlisting}[caption={Tire force relaxation via EMA}]
const relaxationLength = params.tireRelaxationLength || 0.3;
const vLongAbs = Math.max(Math.abs(wheelLongitudinalSpeed), 0.5);
const tireRelaxTau = Math.max(relaxationLength / vLongAbs, 1e-4);
const tireRelaxAlpha = clamp01(1.0 - Math.exp(-dt / tireRelaxTau));

const prev = state.prevTireForce[name];
const relaxedLat = prev.lat + (scaledLateral - prev.lat) * tireRelaxAlpha;
const relaxedLon = prev.lon + (scaledLongitudinal - prev.lon) * tireRelaxAlpha;
prev.lat = relaxedLat;
prev.lon = relaxedLon;
\end{lstlisting}

\textbf{Notes}:
\begin{itemize}
  \item Relaxation is implemented via exponential moving average (EMA) filtering
  \item The EMA blending factor is computed from the time constant
  \item Relaxation is applied AFTER the Pacejka formula but BEFORE output
  \item This models the real phenomenon where tires lag when slip suddenly changes
  \item At very low speeds ($< 0.5$ m/s), relaxation is still applied with a minimum tau to prevent stalling
\end{itemize}

\subsection{Grip State Machine}

\label{sec:grip-state-machine}

The tire system maintains a three-state automaton tracking grip state based on friction utilization. See Section \ref{sec:hysteretic-grip-state} for detailed state equations and transitions.

\textbf{States}:
\begin{itemize}
  \item \textbf{stable}: Wheel has grip available, $u < 0.65$
  \item \textbf{slipping}: Wheel is losing/has lost grip, $u \ge 0.90$
  \item \textbf{recovering}: Wheel is regaining grip, $0.80 \le u < 0.90$ (from slipping state)
\end{itemize}

\textbf{Thresholds} (hysteresis prevents oscillation):
\begin{itemize}
  \item Enter slipping: $u > 0.90$
  \item Exit slipping (enter recovering): $u < 0.80$
  \item Exit recovering (enter stable): $u < 0.65$
\end{itemize}

\subsection{Per-Wheel Grip Smoothing}

The friction circle utilization is smoothed with a state-dependent EMA coefficient:

\begin{equation}
  u_{\text{smooth}}(t) = u_{\text{smooth}}(t-\Delta t) + (u(t) - u_{\text{smooth}}(t-\Delta t)) \cdot \alpha_{\text{grip}}(s)
  \label{eq:grip-smoothing}
\end{equation}

\textbf{EMA Coefficients}:
\begin{equation}
  \alpha_{\text{grip}}(s) = \begin{cases}
    0.15 & \text{if } s = \text{stable} \\
    0.35 & \text{if } s = \text{slipping or recovering}
  \end{cases}
  \label{eq:grip-ema-rate}
\end{equation}

\textbf{Physical Meaning}: During stable grip, smooth slowly (preserve detail). During slipping, smooth faster (reduce noise from Pacejka nonlinearity). This adaptive filtering significantly improves driving feel.

\textbf{Source Code}: \coderef{tires.js}{374--377}

\textbf{Notes}:
\begin{itemize}
  \item The slow stable rate (0.15) allows fine control and feeling of grip transitions
  \item The fast slipping rate (0.35) smooths constraint solver noise without lag
  \item These rates are independent parameters, tunable via UI
  \item The smoothing affects visual feedback (HUD indicators) but not physics calculations
\end{itemize}

\subsection{Low-Speed Lateral Force Fade}

\label{sec:lateral-fade}

At very low speeds ($< 2$ m/s), lateral forces are attenuated to prevent twitching from constraint solver noise:

\begin{equation}
  \text{fade}(v) = \max\left(0, \min\left(1, \frac{v}{2.0}\right)\right)
  \label{eq:lateral-fade-function}
\end{equation}

\begin{equation}
  F_{\text{lat,final}} = F_{\text{lat}} \cdot \text{fade}(|v|)
  \label{eq:lateral-fade-applied}
\end{equation}

\textbf{Source Code}: \coderef{tires.js}{269--276}

\textbf{Parameters}:
\begin{itemize}
  \item Fade threshold = 2.0 m/s
  \item Below 2 m/s: lateral force smoothly decreases to zero
  \item Above 2 m/s: full lateral force applied
\end{itemize}

\textbf{Notes}:
\begin{itemize}
  \item This is a practical solution to the slip angle denominator clamp (Section \ref{eq:slip-angle})
  \item At standstill, the constraint solver produces large micro-motions
  \item These micro-motions create large slip angles from the denominator clamp
  \item Fading lateral force prevents the car from twitching when parked
\end{itemize}

\newpage
